\chapter{Análisis exploratorio del corpus}
\label{anx:eda}

Este anexo recoge el detalle del análisis exploratorio del Financial PhraseBank
resumido en el capítulo~\ref{cap:datos}. Las estadísticas del corpus proceden de
\texttt{reports/results/eda\_corpus.json}, y el reparto de clases sobre las
frases finalmente utilizadas se obtiene de los tres ficheros de partición
generados a partir de él.

\begin{table}[htbp]
  \centering
  \caption{Composición de la variante empleada (acuerdo unánime entre anotadores).}
  \label{tab:anx-corpus}
  \begin{tabular}{lr}
    \toprule
    Concepto & Valor \\
    \midrule
    Frases totales del corpus       & 2.264 \\
    Duplicados exactos eliminados   & 5 \\
    Frases únicas utilizadas        & 2.259 \\
    Registros nulos o vacíos        & 0 \\
    \midrule
    \multicolumn{2}{l}{\textit{Reparto de clases sobre las 2.259 frases utilizadas}} \\
    Frases negativas                & 303 (13,4\%) \\
    Frases neutras                  & 1.386 (61,4\%) \\
    Frases positivas                & 570 (25,2\%) \\
    \bottomrule
  \end{tabular}
\end{table}

Los cinco duplicados exactos eliminados pertenecían todos a la clase neutra, de
modo que su retirada no altera de forma apreciable las proporciones del corpus.

\begin{table}[htbp]
  \centering
  \caption{Longitud de las frases del corpus.}
  \label{tab:anx-longitudes}
  \begin{tabular}{lrrrrr}
    \toprule
    Medida & Media & Mediana & Desv. típica & Mínimo & Máximo \\
    \midrule
    Palabras    & 22,4  & 21,0  & 10,1 & 2 & 81  \\
    Caracteres  & 122,0 & 111,0 & 56,2 & 9 & 315 \\
    \bottomrule
  \end{tabular}
\end{table}

El corpus se distribuye en cuatro variantes según el grado de acuerdo exigido
entre anotadores. La tabla~\ref{tab:anx-variantes} las compara. Se eligió la más
exigente por ser la menos ambigua, asumiendo que es también la más pequeña.

\begin{table}[htbp]
  \centering
  \caption{Comparación de las cuatro variantes del Financial PhraseBank. Los
           recuentos por clase corresponden al total de cada variante, antes de
           eliminar duplicados exactos.}
  \label{tab:anx-variantes}
  \begin{tabular}{lrrrrr}
    \toprule
    Variante & Totales & Únicas & Negativas & Neutras & Positivas \\
    \midrule
    Acuerdo unánime (100\%) & 2.264 & 2.259 & 303 & 1.391 & 570 \\
    Acuerdo del 75\%        & 3.453 & 3.448 & 420 & 2.146 & 887 \\
    Acuerdo del 66\%        & 4.217 & 4.211 & 514 & 2.535 & 1.168 \\
    Acuerdo del 50\%        & 4.846 & 4.838 & 604 & 2.879 & 1.363 \\
    \bottomrule
  \end{tabular}
\end{table}

\begin{figure}[htbp]
  \centering
  \includegraphics[width=0.8\textwidth]{class_distribution.png}
  \caption{Distribución de clases en la variante de acuerdo unánime. El
           desequilibrio en contra de la clase negativa es el que obliga a
           ponderar las clases durante el entrenamiento.}
  \label{fig:anx-clases}
\end{figure}

\begin{figure}[htbp]
  \centering
  \includegraphics[width=0.8\textwidth]{text_length_distribution.png}
  \caption{Distribución de la longitud de las frases, próxima al formato de un
           titular.}
  \label{fig:anx-longitud}
\end{figure}

\begin{figure}[htbp]
  \centering
  \includegraphics[width=0.8\textwidth]{text_length_boxplot.png}
  \caption{Longitud de las frases por clase. Las tres clases presentan
           longitudes comparables, de modo que la longitud no actúa como pista
           espuria para el clasificador.}
  \label{fig:anx-boxplot}
\end{figure}

\begin{figure}[htbp]
  \centering
  \includegraphics[width=0.8\textwidth]{config_comparison.png}
  \caption{Tamaño y reparto de clases de las cuatro variantes del corpus.}
  \label{fig:anx-variantes-fig}
\end{figure}

\chapter{Detalle del núcleo de modelado}
\label{anx:modelado}

Este anexo completa el capítulo~\ref{cap:modelado} con la configuración de los
modelos y con el desglose por clases que allí se resume. Las cifras proceden de
\texttt{reports/results/baseline.json}, \texttt{transformer.json} y
\texttt{llm\_zeroshot.json}.

\section{Hiperparámetros del entrenamiento}
\label{anx:hiperparametros}

\begin{table}[htbp]
  \centering
  \caption{Configuración del ajuste fino de DeBERTa-v3-base.}
  \label{tab:anx-hiper}
  \begin{tabularx}{\textwidth}{lX}
    \toprule
    Parámetro & Valor \\
    \midrule
    Modelo de partida           & \texttt{microsoft/deberta-v3-base} \\
    Longitud máxima de entrada  & 128 unidades léxicas \\
    Épocas                      & 10 \\
    Tasa de aprendizaje         & 0,00002 \\
    Tamaño de lote              & 16 \\
    Acumulación de gradiente    & 1 \\
    Proporción de calentamiento & 0,1 \\
    Decaimiento de pesos        & 0,01 \\
    Precisión numérica          & completa (32 bits; media precisión desactivada) \\
    Peso de la clase negativa   & 2,4858 \\
    Peso de la clase neutra     & 0,5433 \\
    Peso de la clase positiva   & 1,3208 \\
    Criterio de selección       & macro-F1 sobre validación \\
    Tiempo de entrenamiento     & 3 minutos en GPU \\
    \bottomrule
  \end{tabularx}
\end{table}

La precisión numérica completa no es la opción habitual, ya que la media
precisión acelera el entrenamiento, pero resultaba incompatible con la función
de pérdida ponderada empleada para compensar el desequilibrio de clases.

\begin{table}[htbp]
  \centering
  \caption{Configuración del modelo base y del modelo de lenguaje sin
           entrenamiento.}
  \label{tab:anx-hiper-otros}
  \begin{tabularx}{\textwidth}{lX}
    \toprule
    Parámetro & Valor \\
    \midrule
    \multicolumn{2}{l}{\textit{Modelo base: TF-IDF y LinearSVC}} \\
    Rango de n-gramas            & 1 a 2 \\
    Frecuencia mínima de término & 2 \\
    Escalado sublineal           & activado \\
    Tamaño del vocabulario       & 5.208 términos \\
    Parámetro de regularización  & 0,5 \\
    Alternativa descartada       & regresión logística (macro-F1 de 0,8547) \\
    \midrule
    \multicolumn{2}{l}{\textit{Modelo de lenguaje sin entrenamiento}} \\
    Modelo                       & \texttt{gpt-4o-mini} \\
    Temperatura                  & 0 \\
    Frases evaluadas             & 339 \\
    Respuestas no interpretables & 0 \\
    Coste total                  & 0,0108 dólares \\
    Tiempo total                 & 266,8 segundos \\
    \bottomrule
  \end{tabularx}
\end{table}

\section{Métricas por clase}
\label{anx:por-clase}

La tabla~\ref{tab:anx-porclase} recoge el desglose completo sobre el conjunto de
prueba, 339 frases que ningún modelo utilizó durante su desarrollo. La columna de
apoyo indica cuántas frases de cada clase contiene ese conjunto.

\begin{table}[htbp]
  \centering
  \caption{Precisión, cobertura y F1 por clase de los tres enfoques sobre el
           conjunto de prueba.}
  \label{tab:anx-porclase}
  \begin{tabular}{llrrrr}
    \toprule
    Enfoque & Clase & Precisión & Cobertura & F1 & Apoyo \\
    \midrule
    \multirow{3}{*}{TF-IDF y LinearSVC}
      & Negativa & 0,8750 & 0,7609 & 0,8140 & 46 \\
      & Neutra   & 0,9050 & 0,9615 & 0,9324 & 208 \\
      & Positiva & 0,8590 & 0,7882 & 0,8221 & 85 \\
    \midrule
    \multirow{3}{*}{DeBERTa-v3-base ajustado}
      & Negativa & 0,9583 & 1,0000 & 0,9787 & 46 \\
      & Neutra   & 0,9952 & 0,9904 & 0,9928 & 208 \\
      & Positiva & 0,9762 & 0,9647 & 0,9704 & 85 \\
    \midrule
    \multirow{3}{*}{gpt-4o-mini sin entrenamiento}
      & Negativa & 0,9583 & 1,0000 & 0,9787 & 46 \\
      & Neutra   & 0,9623 & 0,9808 & 0,9714 & 208 \\
      & Positiva & 0,9747 & 0,9059 & 0,9390 & 85 \\
    \bottomrule
  \end{tabular}
\end{table}

Dos lecturas merecen destacarse. La primera es que el modelo base recupera solo
el 76,09\% de las frases negativas, mientras que los otros dos las identifican
todas. La segunda es que el modelo de lenguaje sin entrenamiento coincide con el
ajustado en la clase negativa, cifra por cifra, y su diferencia se concentra por
completo en la clase positiva, donde su cobertura baja al 90,59\%.

\section{Matrices de confusión}
\label{anx:confusion}

\begin{figure}[htbp]
  \centering
  \includegraphics[width=0.8\textwidth]{baseline_confusion_matrix_linearsvc.png}
  \caption{Matriz de confusión del modelo base sobre el conjunto de prueba.}
  \label{fig:anx-cm-baseline}
\end{figure}

\begin{figure}[htbp]
  \centering
  \includegraphics[width=0.8\textwidth]{transformer_confusion_matrix_deberta.png}
  \caption{Matriz de confusión de DeBERTa-v3-base ajustado. La fila de la clase
           negativa no presenta ningún fallo.}
  \label{fig:anx-cm-deberta}
\end{figure}

\begin{figure}[htbp]
  \centering
  \includegraphics[width=0.8\textwidth]{llm_zeroshot_confusion_matrix.png}
  \caption{Matriz de confusión de gpt-4o-mini sin entrenamiento. Los errores se
           concentran en frases positivas clasificadas como neutras.}
  \label{fig:anx-cm-llm}
\end{figure}

\chapter{Contraste empírico: figuras}
\label{anx:backtest}

Este anexo recoge las figuras de los tres análisis del capítulo~\ref{cap:backtest},
calculados sobre la ventana comprendida entre el 25 de junio de 2025 y el 24 de
junio de 2026.

\begin{figure}[htbp]
  \centering
  \includegraphics[width=0.8\textwidth]{backtest_leadlag.png}
  \caption{Correlación entre la señal de sentimiento y el rendimiento del precio
           a distintos desfases. Los desfases negativos comparan con días
           anteriores y los positivos con días posteriores. La relación
           detectable se sitúa a la izquierda del cero, es decir, hacia el
           pasado.}
  \label{fig:anx-leadlag}
\end{figure}

\begin{figure}[htbp]
  \centering
  \includegraphics[width=0.8\textwidth]{backtest_eventstudy.png}
  \caption{Diferencia de rendimiento tras las jornadas de sentimiento más bajo,
           a uno, tres y cinco días, por empresa y en conjunto. Ninguna
           diferencia alcanza significación estadística.}
  \label{fig:anx-eventstudy}
\end{figure}

\begin{figure}[htbp]
  \centering
  \includegraphics[width=0.8\textwidth]{backtest_casestudy.png}
  \caption{Caso de estudio visual sobre dos valores de la cartera, con la
           evolución conjunta de la señal de sentimiento y del precio.}
  \label{fig:anx-casestudy}
\end{figure}

\chapter{Despliegue}
\label{anx:despliegue}

Este anexo detalla el empaquetado y los puntos de acceso del servidor descritos
en el capítulo~\ref{cap:app}.

\section{Empaquetado}
\label{anx:empaquetado}

La imagen parte de \texttt{python:3.12-slim} y se construye en capas ordenadas de
menor a mayor frecuencia de cambio, de modo que una modificación del código no
obligue a reinstalar las dependencias.

La primera capa instala PyTorch desde el índice de distribución específico para
procesadores convencionales, en lugar de la versión con soporte para tarjetas
gráficas que se obtendría por defecto. En explotación no se entrena nada, y el
trabajo pesado de clasificación y de extracción pertenece a la fase de
preparación de los datos, no al servicio en funcionamiento, de modo que la
variante para procesadores convencionales resulta suficiente y reduce el tamaño
final de la imagen de forma apreciable. La segunda capa instala el resto de
dependencias de ejecución desde un fichero de requisitos propio de producción,
del que se han excluido las herramientas de análisis y experimentación. Se
instala además la biblioteca del sistema que necesita el índice vectorial para
paralelizar.

La tercera capa copia el código fuente y la cuarta incorpora la base de datos
SQLite dentro de la propia imagen, lo que permite que el contenedor arranque con
todos sus datos sin depender de almacenamiento externo. La quinta descarga el
modelo de clasificación desde Hugging Face Hub durante la construcción y lo deja
en la caché local, evitando así una descarga de varios cientos de megabytes en
cada reactivación del servicio.

El contenedor se ejecuta con un usuario sin privilegios de administrador y
escucha en el puerto 7860 por defecto, valor que puede sustituirse mediante una
variable de entorno, que es el mecanismo por el que la plataforma de despliegue
asigna el suyo.

\section{Puntos de acceso de la API}
\label{anx:endpoints}

El servidor expone cinco puntos de acceso de datos y sirve además la interfaz
web desde la raíz.

\begin{table}[htbp]
  \centering
  \caption{Puntos de acceso del servidor.}
  \label{tab:anx-endpoints}
  \begin{tabularx}{\textwidth}{lX}
    \toprule
    Ruta & Función \\
    \midrule
    \texttt{/health} & Comprueba que el servicio responde y devuelve el entorno activo. \\
    \texttt{/catalog} & Lista las empresas disponibles con su volumen y rango de fechas. \\
    \texttt{/portfolio} & Resumen por empresa y bandera de alerta para la ventana elegida. \\
    \texttt{/company/\{ticker\}} & Serie de sentimiento, noticias recientes y eventos de una empresa. \\
    \texttt{/company/\{ticker\}/briefing} & Devuelve el informe pre-generado de esa empresa. \\
    \texttt{/} & Sirve la interfaz web. \\
    \bottomrule
  \end{tabularx}
\end{table}

El catálogo es dinámico: cualquier empresa con noticias relevantes en la base de
datos aparece en él sin necesidad de modificar el código. El punto de acceso de
la cartera admite dos parámetros, uno para fijar la ventana de observación y
otro para seleccionar el subconjunto de empresas que se desea consultar. Los
ficheros estáticos de la interfaz se sirven desde la ruta \texttt{/static}.

\section{Repositorio y modelo}
\label{anx:repositorio}

El código completo del proyecto, incluidos los cuadernos de análisis, los
scripts de ingesta y la aplicación, está disponible en el repositorio público
\texttt{github.com/Rhulli34/TFM}.

El clasificador de sentimiento ajustado en el capítulo~\ref{cap:modelado} está
publicado como modelo público en Hugging Face Hub, bajo el identificador
\texttt{Rhulli/financial-news-radar-deberta}, y puede descargarse y reutilizarse
de forma independiente al resto del sistema.
